% Round 12, numerical-lift branch.  New files only; older rounds are unchanged.
% This is an internal proof candidate.  The analytic conclusion explicitly
% depends on the Round 10/11 actual-form foundations, not on a published G theorem.
\section{Curve lifts on regular SNC models, with multiplicities}
\label{r12n:section}

\paragraph{Precise gain and status.}
The reduced semistable curve correction was already an internal result
of Round 9.  We prove its multiplicity-weighted extension below and then
construct the original-model virtual vector bundle required by (OC4).
This supplies a new connection to the actual Green-form problem rather
than another formulation of (OC4).  In relative dimension three it
applies to every geometrically homologically trivial codimension-two
cycle on a supplied regular projective SNC model.  The word ``SNC''
refers to the reduced support; the scheme-theoretic fibre may have
arbitrary positive multiplicities.  No tameness assumption is imposed.
The algebraic argument uses published divisor, intersection, and trace
theorems.  The passage from the constructed witness to actual forms
uses the expressly identified internal Round 10/11 results.

Let $K/\mathbb Q_p$ be finite, let $R=K^\circ$, and let $k$ be its
finite residue field.  Let $\mathcal X/R$ be a supplied integral,
regular, flat, projective model of a smooth projective geometrically
integral $n$-fold $X/K$, where $n\geq3$.  Assume
\[
       \mathcal X_k=\sum_{i\in I}m_iD_i,\qquad m_i\in\mathbb Z_{>0},
       \qquad Y=(\mathcal X_k)_{\mathrm{red}}=\bigcup_iD_i
\tag{NL1}\label{r12n:fibre}
\]
has strict normal crossings and all its closed strata are smooth over
$k$.  Write $j_i:D_i\hookrightarrow\mathcal X$.  All Chow and numerical
groups have rational coefficients.  Fix one relatively ample line
bundle $H$ on $\mathcal X$ and use its restrictions on every stratum.
All degrees are native degrees over $k$.

\begin{theorem}[Original-model curve witness and strong G]
\label{r12n:main}
Let $Z$ be a prescribed rational codimension-$(n-1)$ cycle on $X$,
and suppose, for one $\ell\ne p$, that
\[
          \operatorname{cl}_{\ell}(Z)=0
          \quad\hbox{in }H^{2n-2}(X_{\bar K},\mathbb Q_\ell(n-1)).
\]
Then there is a rational virtual vector bundle
$\xi\in K_0(\mathcal X)_{\mathbb Q}$ such that
\[
 \operatorname{ch}_{n-1}(\xi|_X)=[Z],\qquad
 \deg\bigl(\operatorname{ch}_{n-1}(\xi|_V)\cap[V]\bigr)=0
\tag{NL2}\label{r12n:witness}
\]
for every integral $(n-1)$-dimensional closed subvariety
$V\subset\mathcal X_k$.  Consequently, under the actual-form
foundations established internally in Rounds 10 and 11,
\[
       \mathfrak o_X^{n-1}([Z])=0,
\]
and strong G holds on exactly $X^{\mathrm{an}}\setminus|Z|^{\mathrm{an}}$
over the original $K$, with its identity tested against all actual
compactly supported test forms.
\end{theorem}

The proof occupies the following subsections.  It constructs the
witness on $\mathcal X$ itself.  There is no analytic pushforward from
an alteration in this proof, and the theorem does not assert that every
smooth $X/K$ admits a regular SNC model over $R$.

\subsection{The weighted divisor kernel}

Put $V^a=\bigoplus_iN^a(D_i)_{\mathbb Q}$ and define the actual
vertical incidence operator
\[
 A_a:V^a\longrightarrow V^{a+1},\qquad
       A_a(v)=\left(j_i^!\sum_jj_{j*}v_j\right)_i.
\tag{NL3}\label{r12n:actual-A}
\]
This notation is legitimate on numerical groups: restriction and
Gysin maps between the smooth strata preserve numerical equivalence
by projection, and the diagonal entry is the first Chern class of
$N_{D_i/\mathcal X}$.  These descriptions also define $A_a$ directly.
The principal-fibre relation in the regular model gives
\[
 m_i c_1(N_{D_i/\mathcal X})
            =-\sum_{j\ne i}m_j[D_i\cap D_j].
\tag{NL4}\label{r12n:normal}
\]

For each connected component $D_e$ of a double intersection, orient
$e=(i,j)$ by $i<j$.  Likewise let $D_\tau$ denote a connected triple
stratum.  Write $C^q$ for the rational space of functions on the
connected $(q+1)$-fold strata, with the usual signed maps
$d_0:C^0\to C^1$, $d_1:C^1\to C^2$.  In particular
$(d_0c)_{ij}=c_j-c_i$ and
$(d_1h)_{ijk}=h_{ij}-h_{ik}+h_{jk}$, with the indicated edges restricted
to the connected triple in question.  Define
\[
 (\delta v)_e=v_j|_{D_e}-v_i|_{D_e},\qquad
       \delta:V^1\longrightarrow\bigoplus_eN^1(D_e).
\]
Let $\gamma_0:C^1\to V^1$ be signed Gysin on constants: an edge
has coefficient $-[D_e]$ on $D_i$ and $+[D_e]$ on $D_j$.
On actual cycles, not merely on numerical groups,
\[
                 \sum_i j_{i*}(\gamma_0h)_i=0.
\tag{NL5}\label{r12n:Jgamma}
\]
Indeed both terms of each edge are the same closed intersection
cycle in $\mathcal X$, with opposite signs.

Let $M$ be multiplication by $m_i$ on the component summands and set
\[
 \widetilde A=M A_1 M,
 \qquad (\widetilde A w)_i
        =\sum_{j\ne i}m_im_j[D_i\cap D_j](w_j-w_i).
\tag{NL6}\label{r12n:weighted-A}
\]
Here products with an intersection mean restriction followed by Gysin.
If $D_c$ on $C^1$ has diagonal entry $m_im_j$ at an edge $(i,j)$, put
\[
             \gamma_0^m=M^{-1}\gamma_0D_c.
\tag{NL7}\label{r12n:weighted-gamma}
\]
Thus an edge constant contributes $-m_j[D_e]h_e$ on $D_i$ and
$+m_i[D_e]h_e$ on $D_j$.  Equation \eqref{r12n:Jgamma} implies
$\widetilde A\gamma_0^m=0$.

\begin{lemma}[Multiplicity-corrected divisor kernel]
\label{r12n:kernel}
There are equalities of rational numerical subspaces
\[
 \ker\widetilde A=\ker\delta+\operatorname{im}\gamma_0^m,
 \qquad
 \boxed{\ \ker A_1=M(\ker\delta)+\operatorname{im}\gamma_0.\ }
\tag{NL8}\label{r12n:kernel-formula}
\]
\end{lemma}

\begin{proof}
We give all weights, since simply reusing the reduced-fibre formula
would give the wrong compatible summand.  On a double stratum $D_{ij}$,
restriction of the principal-fibre relation gives
\[
 m_jN_{D_{ij}/D_i}+m_iN_{D_{ij}/D_j}
                  +\sum_{k\ne i,j}m_k[D_{ijk}]=0
                  \quad\hbox{in }N^1(D_{ij}).
\tag{NL9}\label{r12n:double-normal}
\]
For an ordered triple $i<j<k$, its contribution to
$\delta\gamma_0^m h$ on the edges $ij,ik,jk$ is, respectively,
\[
        -m_k(d_1h)_{ijk}[D_{ijk}],\quad
        +m_j(d_1h)_{ijk}[D_{ijk}],\quad
        -m_i(d_1h)_{ijk}[D_{ijk}].
\tag{NL10}\label{r12n:triple-signs}
\]
For example, the coefficient of $h_{ij}$ on edge $ij$ is the sum
$m_iN_{D_{ij}/D_j}+m_jN_{D_{ij}/D_i}$; the other two edge constants
give $+m_kh_{ik}$ and $-m_kh_{jk}$ there.  This proves the displayed
signs and shows that $d_1h=0$ implies $\delta\gamma_0^m h=0$.

Define positive rational weights
\[
 a_e=m_im_j\deg_{D_e}H^{n-1},\qquad
 b_\tau=m_im_jm_k\deg_{D_\tau}H^{n-2}
\]
and let $D_a,D_b$ be the corresponding diagonal matrices.  Set
\[
 \lambda(w)_e=m_im_j\deg_{D_e}((\delta w)_eH^{n-2}).
\]
The projection formula and \eqref{r12n:triple-signs} give
\[
 \lambda(w)^th
       =\sum_i m_i\deg_{D_i}\bigl(w_i(\gamma_0^m h)_iH^{n-2}\bigr),
 \qquad \lambda\gamma_0^m=-d_1^tD_bd_1.
\tag{NL11}\label{r12n:lambda}
\]
Also
\[
 \deg_{D_i}((\gamma_0^m h)_iH^{n-1})
                  =m_i^{-1}(d_0^tD_ah)_i,
 \quad
 (d_0^t\lambda(w))_i
                  =-\deg_{D_i}((\widetilde Aw)_iH^{n-2}).
\tag{NL12}\label{r12n:divergence}
\]

Take $h\in\ker d_1\cap\ker(d_0^tD_a)$ and write
$u=\gamma_0^m h$.  The triple relation gives $\delta u=0$, and
\eqref{r12n:divergence} makes each $u_i$ primitive with respect to
$H$.  By \eqref{r12n:lambda},
\[
             \sum_i m_i\deg_{D_i}(u_i^2H^{n-2})
                       =h^t\lambda(u)=0.
\]
The divisor Hodge index theorem is negative definite on the primitive
numerical divisor space of each smooth projective $D_i$.  Since
$m_i>0$, every $u_i$ is zero.  Thus
\[
       \gamma_0^m\bigl(\ker d_1\cap\ker(d_0^tD_a)\bigr)=0.
\tag{NL13}\label{r12n:harmonic-zero}
\]

Now let $w\in\ker\widetilde A$.  By
\eqref{r12n:divergence}, $\lambda(w)$ annihilates $\operatorname{im}d_0$.
The weighted finite-dimensional decomposition
\[
 \ker d_1=\operatorname{im}d_0\ \oplus\
           \bigl(\ker d_1\cap\ker(d_0^tD_a)\bigr)
\]
and \eqref{r12n:lambda}--\eqref{r12n:harmonic-zero} show that it
annihilates all of $\ker d_1$.  Positivity of $D_b$ gives
\[
 \lambda(w)\in\operatorname{im}d_1^t
                      =\operatorname{im}(d_1^tD_bd_1).
\]
The equality follows from equality of kernels over $\mathbb R$ and
rational rank; a rational right-hand side has a rational solution.
Choose $a\in C^1$ with $d_1^tD_bd_1a=\lambda(w)$, and put
$w_0=w+\gamma_0^m a$.  Then
$\widetilde Aw_0=0$ and $\lambda(w_0)=0$.

Consequently $(\delta w_0)_e$ is primitive on each smooth projective
$(n-1)$-fold $D_e$.  Directly summing the two endpoint contributions
in \eqref{r12n:weighted-A} yields
\[
 \begin{split}
 &\sum_{e=(i,j)}m_im_j
       \deg_{D_e}((\delta w_0)_e^2H^{n-3})\\
 &\hspace{12mm}=-\sum_i
       \deg_{D_i}\bigl(w_{0,i}(\widetilde Aw_0)_iH^{n-3}\bigr)=0.
 \end{split}
\tag{NL14}\label{r12n:second-energy}
\]
Divisor Hodge index on the double strata and positivity of $m_im_j$
now give $\delta w_0=0$.  This proves
$\ker\widetilde A\subset\ker\delta+\operatorname{im}\gamma_0^m$.
The reverse inclusion follows from \eqref{r12n:weighted-A} and
$\widetilde A\gamma_0^m=0$.  Finally multiply by $M$ and use
$M\gamma_0^m=\gamma_0D_c$, whose image equals
$\operatorname{im}\gamma_0$ because $D_c$ is invertible.
\end{proof}

\paragraph{Known positivity and native fields.}
The only positivity used is the divisor Hodge index theorem, first on
$D_i$ and then on $D_e$.  For $n=3$, the exponent $n-3$ in
\eqref{r12n:second-energy} is zero.  There is no appeal to a
standard-conjecture sign statement for middle-dimensional cycles.
A primary reference is Kleiman, \emph{The Picard scheme}, Theorem
6.3(a,b,h) and Theorem B.27.  The surface index theorem, followed by
a general complete-intersection surface over $\bar k$, proves
nonpositivity in higher dimension; 6.3(h) identifies the zero locus
with numerical triviality.  Native numerical equivalence can be
checked after finite field extension by norming a test cycle, so the
same conclusion holds over $k$ without assuming geometrically
connected strata.  Our connected-stratum spaces and all their
weights were defined over $k$; no normalization by the number of
geometric components has been made.

\subsection{Compatible divisor classes and the multiplicity-weighted trace}

\begin{lemma}[Cohomological lift of a numerically compatible tuple]
\label{r12n:cohom-lift}
If $w\in\ker\delta$, its component divisor classes are restrictions
of a class
\[
      b_Y\in H^2(Y_{\bar k},\mathbb Q_\ell(1)).
\]
This class lifts to
$b_{\mathcal X}\in H^2(\mathcal X_{R^{\mathrm{sh}}},\mathbb Q_\ell(1))$.
No algebraicity or Frobenius-fixedness of the chosen lift is required.
\end{lemma}

\begin{proof}
Numerical equivalence of divisors implies equality of their geometric
rational $\ell$-adic classes: by Kleiman 6.3(a,b), a numerically
trivial line bundle has an algebraically trivial nonzero power, and
algebraically trivial divisors have zero rational cycle class.  Thus
the classes of $w_i$ form a compatible tuple.

For the closed cover of $Y$ by its components, the augmented complex
of constant sheaves with terms on the distinct multiple intersections
is exact.  At a geometric point it is the augmented cochain complex
of the simplex of the components containing that point.  Closed
immersion pushforward is exact, so its hypercohomology spectral
sequence is
\[
 E_1^{a,b}=H^b(D^{[a+1]}_{\bar k},\mathbb Q_\ell(1))
             \Longrightarrow H^{a+b}(Y_{\bar k},\mathbb Q_\ell(1)).
\tag{NL15}\label{r12n:cech}
\]
Deligne's purity theorem for smooth projective finite-field varieties
makes $E_1^{a,b}$ pure of weight $b-2$.  The compatible tuple lies in
$E_2^{0,2}$, of weight zero.  For $s\geq2$, an outgoing differential
from $(0,2)$ lands in $(s,3-s)$, of weight $1-s$, and hence is zero;
there is no incoming differential from a negative column.  Therefore
the edge map from total $H^2$ surjects onto this compatible kernel,
which gives $b_Y$.  This argument uses neither exactness of the
Frobenius-fixed-vector functor nor a semisimplicity assertion.

The nilpotent closed immersion $Y\hookrightarrow\mathcal X_k$
does not change the etale topos.  Proper base change and the
acyclicity of a strictly henselian local base identify
\[
 H^2(\mathcal X_{R^{\mathrm{sh}}},\mathbb Q_\ell(1))
       \simeq H^2((\mathcal X_k)_{\bar k},\mathbb Q_\ell(1))
       \simeq H^2(Y_{\bar k},\mathbb Q_\ell(1)).
\]
These statements may first be used with $\mathbb Z/\ell^a$ coefficients
and then passed to the usual rational $\ell$-adic cohomology; the
finite cohomology groups on the proper geometric fibre give the
Mittag--Leffler condition.  This yields $b_{\mathcal X}$.
\end{proof}

\begin{lemma}[Trace with the scheme-theoretic multiplicities]
\label{r12n:weighted-trace}
Let $f:\mathcal X\to\operatorname{Spec}R$.  For a class
$c\in H^{2n}(\mathcal X_{R^{\mathrm{sh}}},\mathbb Q_\ell(n))$,
the structural trace satisfies
\[
 \sum_i m_i\operatorname{Tr}_{(D_i)_{\bar k}}
                    (c|_{(D_i)_{\bar k}})
      =\operatorname{Tr}_{X_{\bar K}}(c|_{X_{\bar K}}).
\tag{NL16}\label{r12n:trace}
\]
On a native disconnected geometric fibre of $D_i$, its trace here
is the sum of the geometric component traces; for a native zero-cycle
it is precisely its degree over $k$.
\end{lemma}

\begin{proof}
The proper flat map $f$ is equidimensional of relative dimension
$n$, so SGA 4 XVIII, Theorem 2.9 applies, with its condition
$(*)_n$ satisfied on all of $\mathcal X$.  Its trace is compatible
with arbitrary base change (Var2) and composition (Var3).  Applied
over $R^{\mathrm{sh}}$, the special and geometric generic traces
of the restriction of the same global class are the same section of
the constant sheaf $\mathbb Q_\ell$.  It remains to identify the
special trace, since the etale topos itself does not see nilpotents.

Remove all component intersections and a further closed subset of
dimension less than $n$.  By SGA 4 XVIII, Lemma 2.1, extension from
the remaining open induces an isomorphism in top compactly supported
cohomology and does not alter the trace.  On a dense open of one
geometric component, its nilpotent thickening has generic length
$m_i$.  Over $\bar k$, choose a separating Noether normalization
and shrink its base.  The reduced open is then finite etale over
an open of $\mathbb A^n$, and the full thickening is finite locally
free.  After an etale base change the reduced cover splits, and each
thickened branch is finite locally free of rank $m_i$.  Such shrinking
is valid by generic flatness; the generic length remains $m_i$.

For this finite locally free map, the unit followed by its trace is
multiplication by $m_i$, by Theorem 2.9, Var4(I).  Since the branch
has the same etale topos as its reduction, this says that its trace
on the corresponding constant sheaf is exactly $m_i$ times the
reduced trace.  Compose with the structural trace of the affine open
using Var3.  The resulting top trace is $m_i$ times the smooth
reduced-component trace.  Add the disjoint dense opens and use
Lemma 2.1 again.  This proves the left side of
\eqref{r12n:trace}, including each native residue degree.
The same argument with torsion coefficients passes to
$\mathbb Q_\ell$; no prime-to-$m_i$ restriction is needed.
\end{proof}

\paragraph{Why this is not a fixed-vector argument.}
The only Frobenius input is that maps between the distinct weights
in \eqref{r12n:cech} vanish.  We choose any cohomological lift
provided by the edge map, lift it over the strict henselization, and
compare the traces of one cup product.  Its generic factor will
vanish geometrically.  An assertion about an absolute generic
cycle class, weight--monodromy, or a Tate-conjecture surjectivity
statement is unnecessary.

\subsection{Vertical correction of an arbitrary curve extension}

\begin{proposition}[Numerically admissible curve extension]
\label{r12n:admissible}
Let $z\in\operatorname{CH}^{n-1}(X)_{\mathbb Q}$ have zero geometric
$\ell$-adic cycle class.  Every extension
$\Theta\in\operatorname{CH}^{n-1}(\mathcal X)_{\mathbb Q}$ can be
corrected by a rational vertical cycle to an extension
$\Theta^{\mathrm{ad}}$ such that
\[
          j_i^!\Theta^{\mathrm{ad}}=0
              \quad\hbox{in }N^{n-1}(D_i)_{\mathbb Q}
              \quad\hbox{for every }i.
\tag{NL17}\label{r12n:component-zero}
\]
\end{proposition}

\begin{proof}
Put $T_i^{\mathrm{Ch}}=j_i^!\Theta\in\operatorname{CH}^{n-1}(D_i)$,
let $T_i$ be its numerical image, and view $T=(T_i)$ as a functional
on $V^1$ by
\[
                   t_\Theta(v)=\sum_i\deg_{D_i}(T_iv_i).
\]
The functional annihilates $\operatorname{im}\gamma_0$ by the
supported identity \eqref{r12n:Jgamma} and the projection formula.
We check that it also annihilates $M\ker\delta$.

For $w\in\ker\delta$, choose $b_{\mathcal X}$ as in
Lemma~\ref{r12n:cohom-lift}; each $w_i$ can be represented by any
rational divisor, and its etale class is independent of that choice
by the numerical-to-homological implication for divisors used there.
On the regular projective model the
rational Chern-character isomorphism (Stacks, Tag 0FEX) supplies
a virtual bundle whose only nonzero operational Chern-character
component caps $[\mathcal X]$ to $\Theta$ in degree $n-1$.
Its etale Chern character gives a class
\[
 a\in H^{2n-2}(\mathcal X_{R^{\mathrm{sh}}},\mathbb Q_\ell(n-1)),
\]
whose restrictions are $\operatorname{cl}(T_i^{\mathrm{Ch}})$ and
$\operatorname{cl}(z)$, by Chern-character functoriality and regular
Gysin compatibility.  Apply Lemma~\ref{r12n:weighted-trace} to
$c=a\cup b_{\mathcal X}$.  It gives
\[
 \begin{split}
 t_\Theta(Mw)
    &=\sum_i m_i\deg_{D_i}(T_iw_i)\\
    &=\operatorname{Tr}_{X_{\bar K}}
        \bigl(\operatorname{cl}(z)\cup b_{\mathcal X}|_{X_{\bar K}}\bigr)
      =0.
 \end{split}
\tag{NL18}\label{r12n:obstruction-zero}
\]
Only the geometric generic class is used in the last equality.
The cycle class here is taken from the displayed Chow representative;
no cycle-class map out of a higher-codimension numerical group is
asserted or needed.

Lemma~\ref{r12n:kernel} implies $t_\Theta|_{\ker A_1}=0$.
The ordinary numerical pairings on the smooth projective components
give perfect finite-dimensional rational dualities
\[
 (V^1)^*=V^{n-1},\qquad (V^2)^*=V^{n-2}.
\]
Projection on each double stratum, and the identical self-normal
term on the diagonal, show that the transpose of $A_1$ is
$A_{n-2}$.  Thus
\[
                T\in(\ker A_1)^\perp=\operatorname{im}A_{n-2}.
\]
Choose a rational tuple $u\in V^{n-2}$ with $A_{n-2}u=T$ and
choose rational Chow representatives for its entries.  Then
\[
                 \Theta^{\mathrm{ad}}
                    =\Theta-\sum_i j_{i*}u_i
\tag{NL19}\label{r12n:correction}
\]
has unchanged generic Chow class and satisfies
\eqref{r12n:component-zero}.
\end{proof}

\subsection{From this Chow lift to an operational virtual-bundle witness}

\begin{lemma}[Pure operational representative on the original regular model]
\label{r12n:K0}
Let $\alpha\in\operatorname{CH}^r(\mathcal X)_{\mathbb Q}$.  There is
$\xi\in K_0(\mathcal X)_{\mathbb Q}$ satisfying
\[
 \operatorname{ch}_r(\xi)\cap[\mathcal X]=\alpha,
 \qquad \operatorname{ch}_q(\xi)=0\text{ in }A^q(\mathcal X)_{\mathbb Q}
                 \quad(q\ne r).
\tag{NL20}\label{r12n:pure-character}
\]
If the component restrictions $j_i^!\alpha$ are numerically zero on
the smooth $D_i$, then
\[
          \deg\bigl(\operatorname{ch}_r(\xi|_V)\cap[V]\bigr)=0
\]
for every integral $r$-dimensional $V\subset\mathcal X_k$.
\end{lemma}

\begin{proof}
The model is regular, finite dimensional, projective over the regular
noetherian base $R$, and has the resolution property.  Stacks,
Lemma 42.58.1, Tag 0FEX, applies: the composite of rational Chern
character with cap by $[\mathcal X]$ is an isomorphism from
$K_0(\operatorname{Vect}\mathcal X)_{\mathbb Q}$ to the total
homological Chow group.  Take the unique preimage of the
pure-codimension class $\alpha$.  The proof of that lemma uses the
Adams operation $\psi^2$ and uniqueness of the preimage to show that
$2^{-r}\psi^2\xi=\xi$.  Since $\psi^2$ multiplies
$\operatorname{ch}_q$ by $2^q$, all $q\ne r$ components vanish as
\emph{operational classes}, giving \eqref{r12n:pure-character}.
This avoids replacing an operational identity on arbitrary supports
by only an equality after cap with the model's fundamental class.

An integral $V\subset\mathcal X_k$ lies in at least one $D_i$.
Regular Gysin commutes with the operational character, so
\[
       j_i^*\operatorname{ch}_r(\xi)\cap[D_i]=j_i^!\alpha.
\]
Because $D_i$ is smooth, its operational class is identified with
this usual Chow class.  Projection therefore identifies the
supported degree on $V$ with
\[
                    \deg_{D_i}\bigl((j_i^!\alpha)\cdot[V]\bigr)=0.
\]
This tests every integral $r$-dimensional special subvariety,
including singular ones and ones lying in several components.
It is not a test only on whole components or on ample powers.
\end{proof}

\begin{proof}[Proof of Theorem~\ref{r12n:main}]
Take the closure of the prescribed cycle as an initial model Chow
extension.  Proposition~\ref{r12n:admissible} produces
$\Theta^{\mathrm{ad}}$.  Apply Lemma~\ref{r12n:K0} with $r=n-1$.
Generic restriction of \eqref{r12n:pure-character} gives the first
condition of \eqref{r12n:witness}; its last assertion gives the
second condition on every required special test.

Equip a finite virtual-bundle presentation of $\xi$ with its model
metrics.  The Round 10 numerical model-Chern detector gives the
strict actual identity
\[
               \operatorname{ch}_{n-1}(\bar\xi)=0
                         \quad\hbox{in }\mathcal B^{n-1,n-1}(X^{\mathrm{an}}).
\tag{NL21}\label{r12n:actual-zero}
\]
Indeed the degree-$n-1$ Chern character is a homogeneous rational
polynomial in model Chern classes; the detector includes virtual
bundles and tests precisely the operational degrees in
\eqref{r12n:witness}.  Round 11's model-curvature class and
numerical-lift corollary (OC4) now give
$\mathfrak o_X^{n-1}([Z])=0$.  Its exact-support prescribed-cycle
Green-pair theorem and global primitive comparison give the stated
strong G form.  A rational equivalence used to represent $[Z]$
does not enlarge the final exceptional set: the invoked Round 11
theorem is explicitly for the prescribed cycle itself and on its
exact complement, with all actual tests.
\end{proof}

\paragraph{What happens to an already chosen curvature.}
If $\Omega_Z$ is the curvature in any original-model Green-pair
construction for $Z$, its generic Chow character agrees with that
of $\xi$.  Round 11's generic-Chow-zero primitive theorem gives a
global actual $(n-2,n-2)$ form $\Psi$ such that
$D\Psi=\Omega_Z-\operatorname{ch}_{n-1}(\bar\xi)=\Omega_Z$.
Subtracting $\Psi$ from the Green form removes curvature without
introducing a new pole.  For a threefold this is a global actual
$(1,1)$ correction to the $(2,2)$ curvature, not a scalar graph
potential.  The algebraic construction of $\xi$ is the missing
input that makes this use of the Round 11 theorem concrete.

\subsection{A geometric check that multiplicities cannot be dropped}

Start with $\mathbb P^3_R$, blow up a rational point of its special
fibre, and then blow up the smooth intersection of its two special
components.  After the first blowup the components are
$D_0=\operatorname{Bl}_p\mathbb P^3$ and $D_1=\mathbb P^3$, meeting
along $\mathbb P^2$.  The second centre is Cartier in either
component, so their strict transforms remain these same varieties.
The new component is
\[
       F=\mathbb P_{\mathbb P^2}
                    (\mathcal O(-1)\oplus\mathcal O(1)),
\]
and the special fibre is $D_0+D_1+2F$.  The strict transforms
$D_0,D_1$ are disjoint, and each meets $F$ in a section.  Locally the
second blowup is the blowup of $(x,y)$ in $xy=\pi$; the chart
$y=xt$ has $\pi=x^2t$, confirming regularity and multiplicities.

Write $E$ for the exceptional divisor on $D_0$, $h$ for the
hyperplane on the base of $F$, and $\zeta=c_1(\mathcal O_F(1))$.
Choose the tautological convention in which $N_{F/\mathcal X}$
has class $-\zeta$.  The two sections on $F$ have classes
$\zeta+h$ and $\zeta-h$.  Their restrictions of $\zeta$ are $h$
and $-h$, respectively; on the matching section of $D_0$ the class
$E$ restricts to $-h$.  Hence the tuple
\[
               w=(E,\ 0,\ -\tfrac12(\zeta+h))
\]
has $\delta w=0$.  But by the actual self-normal entry
$N_{D_0/\mathcal X}=-2E$, the $D_0$ component of $A_1w$ is
\[
                       -2E^2+E^2=-E^2\ne0.
\]
Here $E^3=1$ detects its nonvanishing.  Thus the unweighted assertion
$\ker\delta\subset\ker A_1$ is false even in this elementary
regular SNC model.  In contrast
$Mw=(E,0,-(\zeta+h))$ is the negative of the ordinary edge-Gysin
class for the edge $D_0F$, and \eqref{r12n:Jgamma} gives $A_1Mw=0$.
This checks the location of $M$ in \eqref{r12n:kernel-formula}; the
general proof also accounts for nonempty triple strata through
\eqref{r12n:triple-signs}.

\subsection{Scope, relation to alterations, and primary references}

\begin{enumerate}
\item The algebraic result covers curves, meaning
$\operatorname{CH}^{n-1}$ on an $n$-fold.  In dimension three this is
all codimension two.  For $n\geq4$ it does not prove all of
$\operatorname{CH}^2$: that correction would require a different
middle-degree kernel.  No substitution of divisor Hodge index for
unknown higher-cycle positivity is made.
\item Reduced strict semistability is the special case $m_i=1$.
For arbitrary positive $m_i$, the new weighted kernel and weighted
trace use only geometric homological vanishing.  The Round 9
nonreduced inner-inverse criterion instead assumed an algebraic
inner inverse and absolute generic vanishing; neither assumption
appears in this curve theorem.
\item The regular model in this section is a model of the original
$X/K$.  A separate Round 12 supported algebraic pushdown theorem
can transport numerical witnesses from an alteration onto a
\emph{supplied regular original model}.  Combining that theorem
with the reduced case above and a semistable alteration extends the
conclusion to that regular-model sector.  It does not justify
inventing a regular original model or pushing an arbitrary actual
form through a proper alteration.
\item The analytic inputs are exactly
\path{round11/foundations/r10_model_chern.tex}, especially its
numerical detector and higher-model-Chern extension, and
\path{round11/root/curvature_class_and_next_step.tex}, (OC2)--(OC4),
together with the exact-support and global-primitive constructions
that file cites.  The present proof fills their numerical-witness
existence hypothesis in the stated sector.  All those analytic
arguments remain internal research results subject to independent
expert review.
\item This local strong-G sector does not by itself establish an
unconditional Beilinson--Bloch height pairing in arbitrary
codimension, its canonicity, a comparison with the Round 9 real
height candidate, or global reciprocity.
\end{enumerate}

The following primary sources were opened and their relevant
statements checked in this round:
\begin{itemize}
\item S. Kleiman, \emph{The Picard scheme}, Theorem 6.3(a,b,h),
and Appendix B, Theorem B.27:
\url{https://indico.ictp.it/event/a0255/session/66/contribution/39/material/0/0.pdf}.
\item P. Deligne, \emph{La conjecture de Weil. I},
Publ. Math. IHES 43 (1974), Theorem (1.6), printed page 276:
\url{https://www.numdam.org/item/PMIHES_1974__43__273_0.pdf}.
\item SGA 4, Expose XVIII, Lemma 2.1, Theorem 2.9, particularly
Var2, Var3, and Var4(I); the proof and Remark 2.10.1 also explicitly
record the role of multiplicities:
\url{https://www.normalesup.org/~forgogozo/SGA4/18/18.pdf}.
\item Stacks Project, Tag 03QO (strict-local acyclicity),
Tag 095T (proper base change), and Tag 0FEX (regular-model rational
operational Chern-character isomorphism and its Adams purity argument):
\url{https://stacks.math.columbia.edu/tag/03QO},
\url{https://stacks.math.columbia.edu/tag/095T},
\url{https://stacks.math.columbia.edu/tag/0FEX}.
\end{itemize}

The two-energy identity, the weighted kernel formula, the weighted
obstruction argument, and their application to (OC4) are deductions
written here, not assertions quoted as theorems from these sources.
